\section{Introduction}\label{sec:intro}

Fault injection is the canonical empirical method for evaluating the
resilience of computer systems to transient hardware
faults~\cite{carreira1998xception,hsueh1997faultinjection,giuffrida2013edfi}.
Its scientific function is twofold: to \emph{validate} resilience
claims by demonstrating that a system survives a targeted
perturbation distribution, and to \emph{falsify} them by exhibiting a
perturbation that the system fails to handle. The two are not
symmetric: validation accumulates evidence under a chosen
perturbation family; falsification produces a single counter-example
sufficient to retract a soundness theorem. The latter, properly
understood in the Popperian tradition~\cite{popper1959logic}, is what
advances the science.

The present report concerns the falsification side. We define an
algebraic perturbation operator, RadFI, designed to produce
\emph{controlled, repeatable, and physically meaningful} Single-Event
Upset emulations against an instrumented Linux filesystem. The
operator is the empirical dual of the BEAMFS recovery operator
\(\mathcal{G}\) defined in~\cite{desbrieres2026beamfsv1}: both act on
the same filesystem state space \(\mathcal{S}\); \(\mathcal{G}\)
drives perturbed states back to consistency, \(\varepsilon\) injects
perturbations whose distribution is calibrated to a chosen SEU
cross-section regime.

\subsection{Position of the problem}

A filesystem resilience theorem of the form ``for every
\(s \models\) and every \(e \in \mathcal{E}\), the recovery operator
drives \(e(s)\) to consistency or to fail-closed in finite steps'' is
meaningful only if the perturbation family \(\mathcal{E}\) is
\emph{realizable}. A theorem that quantifies over a family no real
fault can produce is a tautology disguised as a theorem. Conversely,
a theorem whose stated \(\mathcal{E}\) excludes a class of physically
realizable faults is \emph{incomplete}, and the proper response is
to enlarge \(\mathcal{E}\) and revise the proof, not to dismiss the
omitted faults as out of scope.

RadFI is constructed as the empirical instrument that closes this
loop. It produces perturbations \(e \in \mathcal{E}\) that are
simultaneously (i) physically meaningful, in the sense that each
\(e\) corresponds to a single bit-flip on volatile memory in
transit, modelling a canonical
SEU~\cite{dodd2003seu,baumann2005softerrors}; (ii) controllable, in
the sense that the position, timing, and target of each flip are
governed by stated parameters; and (iii) replayable, in the sense
that a fixed PRNG seed reproduces the exact perturbation sequence
across runs, enabling regression on resilience claims.

\subsection{Where we are going}

This report defends a single thesis: a filesystem resilience theorem
must be subject to empirical falsification, and the instrument of
falsification must be a formally specified algebraic operator whose
own faithfulness is theorem-bound. The chain is presented in
\cref{sec:m2c} as a three-tier descent --- mathematics, algorithm,
kernel implementation --- closed by a traceability table where each
line of the perturbation operator's specification is paired with the
algorithmic step and the kprobe-based kernel construct that realizes
it.

\subsection{Contributions}

We claim the following:

\begin{enumerate}
\item A formal definition of the perturbation operator
  \(\varepsilon\) acting on the BEAMFS-compatible state space, with
  stated targeting algebra (region, device, inode, block, op type)
  and a probability sampling regime calibrated to canonical SEU
  rate models (\cref{sec:operator}).
\item A \emph{faithfulness theorem} (\cref{thm:faithfulness}, full
  proof in \cref{app:proof}) showing that under stated hypotheses on
  the bit-flip distribution and the kprobe interception path,
  \(\varepsilon\) is statistically indistinguishable from the
  canonical SEU perturbation operator \(\varepsilon_{\text{SEE}}\)
  of~\cite{desbrieres2026beamfsv1} within bounded probability
  \(\delta\).
\item A three-tier descent (\cref{sec:m2c}) from \(\varepsilon\) to a
  kprobe-based algorithm to a Linux kernel module implementation,
  with a traceability table auditable line by line.
\item An \emph{empirical falsification} (\cref{sec:empirical}) of
  BEAMFS v1 Theorem IV.1~\cite{desbrieres2026beamfsv1} using RadFI
  v0.1.0 against the BEAMFS v1 reference implementation, dated
  2026-04-28; and a \emph{corroboration} of BEAMFS v2 INLINE
  recovery against equivalent disk-level corruption, dated
  2026-04-29.
\end{enumerate}

\subsection{Scope}

This is v1. We model bit-flip perturbations on bio-layer payloads
only. We do not model: bus-level SET on the DRAM channel; SEFI on
storage controller firmware; SEL on the host system; multi-flip
events within a single bio (capped at one flip per intercepted bio
in v0.1.x); RDMA fast-path interception bypassing the bio layer;
cumulative TID/DDD effects (these are below the filesystem and not
perturbations in the SEE sense). RadFI requires kernel privileges to
load and operates on instrumented test devices only; it is not a
remote attack model and is not, by construction, a production
deployment tool.

\subsection{Anti-claim}

We are explicit about what RadFI v1 does \emph{not} establish:

\begin{itemize}
\item \emph{Faithfulness as ground-truth equivalence with beam-time
  measurements.} \cref{thm:faithfulness} shows statistical
  indistinguishability under stated hypotheses, not physical
  equivalence with accelerator-measured upset distributions.
  Bridging this gap requires beam-time at facilities such as RADEF,
  GANIL, or TRIUMF, and is reserved for v2.
\item \emph{Falsification of BEAMFS v2.} The falsification of v1
  Theorem IV.1 in \cref{sec:empirical} does not extend to v2, whose
  revised recovery properties are corroborated empirically but not
  yet formally proved. The successor
  report~\cite{desbrieres2026beamfsv2} states and proves
  Theorem v2.1.
\item \emph{Coverage of adversarial fault injection by an attacker
  with kernel privileges.} RadFI's targeting filters could in
  principle be misused; the v0.1.x build refuses to attach to
  physical block devices and is restricted to loop-mounted test
  images. This is a safety measure, not a security guarantee.
\end{itemize}
